% Part:sets-functions-relations
% Chapter: size-of-sets
% Section: reduction-alt

\documentclass[../../../include/open-logic-section]{subfiles}

\begin{document}

\olfileid[fa-IR]{sfr}{siz}{red-alt}

\olsection{کاهش}

\begin{editorial}
  این بخش ناشمارابودن را از راه کاهش ثابت می‌کند و با نتایجِ
  \olref[nen-alt]{sec} مطابقت دارد. روایتی جایگزین و اندکی مشروح‌تر که
  با نتایجِ \olref[nen]{sec} مطابقت دارد، در \olref[red]{sec} آمده است.
\end{editorial}

با استدلالی قطری ثابت کردیم که $\Bin^\omega$ !!{nonenumerable} است.
با استدلال قطریِ مشابهی نشان دادیم که $\Pow{\Nat}$ !!{nonenumerable}
است. اما راه دیگری نیز برای اثبات اینکه $\Pow{\Nat}$
!!{nonenumerable} است وجود دارد: نشان دهید که \emph{اگر
$\Pow{\Nat}$ !!{enumerable} باشد، آنگاه $\Bin^\omega$ نیز
!!{enumerable} است}. چون می‌دانیم $\Bin^\omega$ !!{nonenumerable} است،
نتیجه می‌شود که $\Pow{\Nat}$ نیز چنین است.

این کار را \emph{کاهش‌دادن} یک مسئله به مسئله‌ای دیگر می‌نامند. در این
مورد، مسئلهٔ برشمردن $\Bin^\omega$ را به مسئلهٔ برشمردن
$\Pow{\Nat}$ کاهش می‌دهیم. راه‌حلی برای مسئلهٔ دوم---یعنی برشماریِ
$\Pow{\Nat}$---راه‌حلی برای مسئلهٔ نخست---یعنی برشماریِ
$\Bin^\omega$---به دست می‌دهد.

برای کاهش مسئلهٔ برشمردن مجموعهٔ~$B$ به مسئلهٔ برشمردن مجموعهٔ~$A$،
روشی ارائه می‌کنیم که برشماریِ~$A$ را به برشماریِ~$B$ تبدیل کند.
آسان‌ترین راه برای این کار تعریف یک !!a{surjection} مانند
$f\colon A \to B$ است. اگر $x_1$، $x_2$، \dots{} مجموعهٔ~$A$ را
برشمارد، آنگاه $f(x_1)$، $f(x_2)$، \dots{} هر عضوِ مجموعهٔ~$B$ را
دست‌کم یک بار در بر خواهد داشت.
\emph{یادداشت دربارهٔ تعریف:} این فهرست ممکن است تکرار داشته باشد و
بنابراین هنوز لزوماً برشماری به معنای تناظر دوسویی نیست. با نگه‌داشتن
نخستین رخدادِ هر مقدار و شماره‌گذاریِ دوباره، برشماریِ دوسوییِ لازم به
دست می‌آید؛ حالتِ تهی جداگانه طبق تعریف پوشش داده می‌شود.
در مورد ما، در پی یک !!a{surjection} مانند $f\colon \Pow{\Nat}
\to \Bin^\omega$ هستیم.

\begin{prob}
نشان دهید اگر یک تابعِ !!{injective} مانند $g\colon B \to A$ وجود داشته
باشد و $B$~!!{nonenumerable} باشد، آنگاه $A$ نیز چنین است. برای این کار
نشان دهید چگونه می‌توان با استفاده از~$g$، برشماریِ~$A$ را به برشماریِ~$B$
تبدیل کرد.
\end{prob}

\begin{proof}[برهانِ {\olref[nen-alt]{thm:nonenum-pownat}} از راه کاهش]
برای انجام کاهش، فرض کنید $\Pow{\Nat}$ !!{enumerable} است و بنابراین
برشماری‌ای از آن به‌صورت $N_{1}$، $N_{2}$، $N_{3}$، \dots وجود دارد.

تابع $f \colon \Pow{\Nat} \to \Bin^\omega$ را با این قرارداد تعریف
کنید که $f(N)$ رشتهٔ $s$ باشد، به‌گونه‌ای که $s(n) = 1$ اگر و تنها
اگر $n \in N$، و در غیر این صورت $s(n) = 0$.
\emph{یادداشت ویراستاری:} زیرنویسِ تعریف‌نشدهٔ متن مبدأ حذف شده است؛
نماد رشته در این تعریف به همان مجموعهٔ ورودی وابسته است.

این امر به‌روشنی تابعی را تعریف می‌کند، زیرا هرگاه $N \subseteq \Nat$،
هر $n \in \Nat$ یا یک !!a{element} از $N$ است یا نیست. برای نمونه،
مجموعهٔ $2\Nat = \Setabs{2n}{n \in \Nat} = \{0,2, 4, 6,
\dots\}$ از اعداد طبیعی زوج به رشتهٔ $1010101\dots$ نگاشته می‌شود؛
$\emptyset$ به $0000\dots$ نگاشته می‌شود و $\Nat$ به
$1111\dots$.

این نگاشت همچنین !!{surjective} است: هر رشته از $0$ها و $1$ها با
مجموعه‌ای از اعداد طبیعی متناظر است؛ یعنی مجموعه‌ای که اعضایش آن اعداد
طبیعی‌اند که با جایگاه‌هایی متناظرند که رشته در آن‌ها $1$ دارد. دقیق‌تر
بگوییم، اگر $s \in \Bin^\omega$، آنگاه $N
\subseteq \Nat$ را چنین تعریف کنید:
\[
N = \Setabs{n \in \Nat}{s(n) = 1}
\]
آنگاه $f(N) = s$؛ این را می‌توان با مراجعه به تعریفِ~$f$ بررسی کرد.

اکنون فهرست زیر را در نظر بگیرید:
\[
f(N_1), f(N_2), f(N_3), \dots
\]
چون $f$ !!{surjective} است، هر عضوِ $\Bin^\omega$ باید برای ورودی‌ای
به‌عنوان مقدارِ~$f$ ظاهر شود و بنابراین باید در فهرست بیاید. پس این
فهرست باید همهٔ~$\Bin^\omega$ را برشمارد.
این نگاشت یک‌به‌یک نیز هست، زیرا رشتهٔ حاصل، عضویتِ هر عدد طبیعی در
مجموعهٔ ورودی را تعیین می‌کند. پس در این مثال، فهرستِ حاصل از یک
برشماریِ دوسویی، خود نیز بدون تکرار است.

بنابراین اگر $\Pow{\Nat}$ !!{enumerable} بود، $\Bin^\omega$ نیز
!!{enumerable} می‌بود. اما $\Bin^\omega$ !!{nonenumerable} است
(\olref[nen-alt]{thm:nonenum-bin-omega}). پس $\Pow{\Nat}$
!!{nonenumerable} است.
\end{proof}

%\begin{explain}
%ممکن است جهتِ کاهش اشتباه گرفته شود.
%برای نمونه، وجود تابعی !!a{surjective} مانند $g \colon \Bin^\omega \to X$
%ناشمارابودنِ $X$ را \emph{ثابت نمی‌کند}. تابعِ $g
%\colon \Bin^\omega \to \Bin$ با ضابطهٔ $g(s) = s(1)$ را در نظر بگیرید؛ این تابع
%هر دنباله از $0$ها و $1$ها را به عضوِ با اندیس یک می‌فرستد.
%این تابع پوشاست، زیرا بعضی دنباله‌ها در این جایگاه $0$ و بعضی
%دیگر $1$ دارند. اما $\Bin$ متناهی است. همچنین تابعِ $f$ باید
%پوشا باشد؛ در غیر این صورت استدلال برقرار نیست:
%$f(x_1)$، $f(x_2)$، \dots{} لزوماً همهٔ
%!!{element}sی~$B$ را در بر نمی‌گیرند. برای نمونه، $h\colon \Nat \to
%\Bin^\omega$ را با ضابطهٔ زیر تعریف کنید:
%\[
%h(n) = \underbrace{000\dots0}_{\text{$n$ بار $0$}}0\ldots
%\]
%این یک تابع است، اما $\Nat$ !!{enumerable} است.
%یادداشت ویراستاری: متنِ غیرفعالِ مبدأ «نخستین» را برای اندیس یک به کار
%برده بود؛ اینجا اندیس صریح بیان شده است. نامِ مجموعهٔ مقصد با زمینه
%هماهنگ و دنبالهٔ خروجی واقعاً نامتناهی شده است. این توضیح همچنان غیرفعال است.
%\end{explain}

\begin{prob}\label{sfr:siz:red-alt:prob:nat-nat}
  با استدلالی کاهشی نشان دهید مجموعهٔ~$X$ از همهٔ توابع
  $f\colon \Nat \to \Nat$ !!{nonenumerable} است. (راهنمایی: یک نگاشت
  پوشا از $X$ به~$\Bin^\omega$ بدهید.)
\end{prob}

\begin{prob}
با استدلالی کاهشی نشان دهید که مجموعهٔ همهٔ \emph{مجموعه‌های} زوج‌های
اعداد طبیعی، یعنی $\Pow{\Nat \times \Nat}$، !!{nonenumerable} است.
\end{prob}

\begin{prob}
با استدلالی کاهشی نشان دهید که $\Nat^\omega$، یعنی مجموعهٔ دنباله‌های
نامتناهیِ اعداد طبیعی، !!{nonenumerable} است.
\end{prob}

%\begin{prob}
%فرض کنید $P$ مجموعهٔ توابع از $\Nat$ به مجموعهٔ $\{0\}$ باشد و $Q$ مجموعهٔ توابعِ \emph{جزئی}
%از اعداد صحیح مثبت به مجموعهٔ $\{0\}$ باشد؛ یعنی این توابع ممکن است برای بعضی
%ورودی‌ها تعریف نشده باشند. نشان دهید $P$~!!{enumerable} است و $Q$ چنین نیست.
%(راهنمایی: مسئلهٔ برشمردن $\Bin^\omega$ را به برشمردنِ~$Q$ کاهش دهید.)
%\end{prob}

\begin{prob}
فرض کنید $S$ مجموعهٔ همهٔ !!{surjection}s از $\Nat$ به مجموعهٔ
$\{0,1\}$ باشد؛ یعنی $S$ از همهٔ !!{surjection}s~$f \colon \Nat
\to \Bin$ تشکیل می‌شود. نشان دهید $S$ !!{nonenumerable} است.
\end{prob}

\begin{prob}
نشان دهید مجموعهٔ~$\Real$ از همهٔ اعداد حقیقی !!{nonenumerable} است.
\end{prob}

\end{document}
